% sec-03: how the program is split between NSF, NIH and private capital.
% Table 18.
\section{How the Program Is Split}

\subsection{The rule that shapes the split}

NSF's SBIR program funds technical innovation with commercial potential, and it
does not fund clinical trials \cite{nsf26510}. The National Institutes of Health
fund clinical research, including small business work that leads to a trial
\cite{nihsbirseed,ncisbir2026}. The program's private capital sits behind a
firewall that gives no investor any say in trial conduct, and pays for what
neither agency pays for \cite{capitalplan2026}. The split in Table 18 follows
those three facts and nothing else. No component is marked yes in two public
columns.

\begin{apptable}
\begin{tabularx}{\textwidth}{>{\raggedright\arraybackslash}p{5.3cm} >{\raggedright\arraybackslash}p{3.2cm} >{\raggedright\arraybackslash}p{3.5cm} >{\raggedright\arraybackslash}X}
\headrow \thead{Component} & \thead{NSF SBIR} & \thead{NIH SBIR} & \thead{Private capital}\\
\midrule
Verification method for the model & Yes, the Pitch & No & No \\
Boundary tests and audit trail & Yes & Cites the NSF result & No \\
Timing tests, hardware in the loop & Yes & No & No \\
Protocol, IND, site activation & No & Yes & In part \\
Drug, surgery, and follow-up & No & Yes, at the site & In part \\
Robot platform, site build-out & No & In part & Yes \\
Commercializing the method & After Phase II & No & Yes \\
\end{tabularx}
\end{apptable}
\vspace{-0.60cm}
\tabcap{Table 18. The program's seven components split between NSF, NIH and private capital,\\
with no component funded twice; the verification work is the part proposed to NSF.}

\subsection{Why the scopes cannot overlap}

The NSF work ends where a patient begins. It produces a verification method, a
test suite and measured failure rates on synthetic cases generated from
published parameters. The NIH work begins where a patient begins: it uses the
verified layer inside a trial that belongs to a site investigator, under an
investigational new drug application filed under 21 CFR Part 312
\cite{ecfr2026part312}. If both were funded, the NIH work would cite the NSF
results rather than repeat them, and each application would disclose the other,
as the SBIR policy directive requires \cite{sbapolicydirective}. Letter 4 tells
NIH so in writing, with NSF copied.

\subsection{The money on each side}

An NSF Phase I is worth up to \$305,000 \cite{nsf26510}. The company's plan
carries an NIH Phase I of \$306,000 and an NIH Phase II of \$1,300,000, for an
SBIR route of \$1,606,000, and places \$5,900,000 of private capital behind the
firewall against the \$2,104,000 delta the route does not buy
\cite{capitalplan2026}. The NSF and NIH Phase I figures are nearly equal and pay
for different work, and neither is counted twice anywhere in the company's
documents. The private figure is a plan, not a completed raise, and the day's
match ledger records its opening balance as zero.

\subsection{What the split means for the trial}

None of the robotic Phase 1 is proposed to NSF. What NSF would fund is the
evidence that the advisory layer inside that trial stays inside its boundaries.
That is the evidence an institutional review board and the FDA will ask for
before the layer is allowed near a patient \cite{fdaaiguidance2025}, and it is
the evidence a clinical funder will ask for before paying for the trial it sits
inside. The NSF work therefore does not compete with the NIH route. It is the
step that makes the NIH route fundable.
